\documentclass[../../main.tex]{subfiles}
\graphicspath{{../../image/}{../../figure/}}
\begin{document}

\section{两个守恒律的真非线性方程组}\label{section:ch17-两律真非线性}

\cref{section:ch17-测度值解与补偿紧致}勾勒的纲领，这里将针对两个守恒律的真非线性方程组\eqref{align:ch17-1731}实现。
特别地，我们的方程组将配备 Riemann 不变量 $(z,w)$ 的坐标系，并按\eqref{align:ch12-riemann归一}规范化，
真非线性条件由\eqref{align:ch12-真非线性}表示，
即 $\lambda_z<0$、$\mu_w>0$。此外，该方程组将配备一个丰富的熵--熵流族，包括\cref{section:ch12-熵熵流对与速矢变换}所构造的 Lax 对，
它们在分析中将起关键作用。

我们将证明：熵条件连同真非线性，可淬灭快速振荡。

\begin{theorem}[两律真非线性的补偿紧致性]\label{theorem:ch17-两律真非线性补偿紧致}
设 $\Omega$ 是 $\mathbb{R}^2$ 的开子集，$\{U_k(x,t)\}$ 是 $L^\infty(\Omega;\mathbb{R}^2)$ 中的有界序列，满足
\begin{align}
\partial_t\eta(U_k)+\partial_x q(U_k)\subset\text{紧集于 }W^{-1,2}_{\mathrm{loc}}(\Omega),
\label{align:ch17-1761}
\end{align}
其中 $(\eta,q)$ 是\eqref{align:ch17-1731}的任意熵--熵流对。% 原书此处误排为 (16.3.1)
则存在子列 $\{U_j\}$，它在 $\Omega$ 上几乎处处收敛。
\end{theorem}

\begin{proof}
应用\cref{theorem:ch17-young测度存在性}，提取子列 $\{U_j\}$ 并识别相联系的 Young 测度族 $\nu_{x,t}$。我们需证明：对几乎所有 $(x,t)$，$\nu_{x,t}$ 的支撑被限制在单点，从而该测度退化为 Dirac 质量。在 Riemann 不变量 $(z,w)$ 平面而非原状态空间上来监测 Young 测度将是方便的。

于是设 $\nu$ 表示任意固定点 $(x,t)\in\Omega$ 处、相对于 $(z,w)$ 变量的 Young 测度，并考虑包含 $\nu$ 支撑的最小矩形 $\mathscr{R}=[z^-,z^+]\times[w^-,w^+]$。我们需证明 $z^-=z^+$ 且 $w^-=w^+$。反证，设 $z^-<z^+$。

我们考虑\eqref{align:ch12-lax族1}的 Lax 熵--熵流对，
此处记为 $(\eta_k,q_k)$，以显式展示对参数 $k$ 的依赖。我们将用 $\eta_k$ 作为权，重新分配 $\nu$ 的质量，把质量向 $\mathscr{R}$ 的边界附近转移。为此，对每个大正整数 $k$，我们在 $\mathscr{R}$ 上联系概率测度 $\nu_k^\pm$，它们通过对连续函数 $h(z,w)$ 的作用定义：
\begin{align}
\langle\nu_k^\pm,h\rangle=\frac{\langle\nu,h\eta_{\pm k}\rangle}{\langle\nu,\eta_{\pm k}\rangle}.
\label{align:ch17-1762}
\end{align}
由于 $\eta_k$ 定义中因子 $e^{kz}$ 的存在，测度 $\nu_k^-$（或 $\nu_k^+$）集中于 $\mathscr{R}$ 的左（或右）侧。当 $k\to\infty$ 时，序列 $\{\nu_k^-\}$ 与 $\{\nu_k^+\}$ 或其子列将在测度空间中弱$^*$ 收敛于概率测度 $\nu^-$ 与 $\nu^+$，它们分别支撑在 $\mathscr{R}$ 的左边 $[z^-]\times[w^-,w^+]$ 与右边 $[z^+]\times[w^-,w^+]$ 上。

对任意固定熵--熵流对 $(\eta,q)$ 与 Lax 对 $(\eta_{\pm k},q_{\pm k})$，应用\eqref{align:ch17-17310}，得
\begin{align}
\langle\nu,q\rangle-\frac{\langle\nu,q_{\pm k}\rangle}{\langle\nu,\eta_{\pm k}\rangle}\langle\nu,\eta\rangle=\frac{\langle\nu,\eta_{\pm k}q-\eta q_{\pm k}\rangle}{\langle\nu,\eta_{\pm k}\rangle}.
\label{align:ch17-1763}
\end{align}
由\eqref{align:ch12-lax族1}与\eqref{align:ch12-递推1}推出
\begin{align}
q_{\pm k}=\bigl[\lambda+O\bigl(\tfrac{1}{k}\bigr)\bigr]\eta_{\pm k}.
\label{align:ch17-1764}
\end{align}

于是令 $k\to\infty$，由\eqref{align:ch17-1763}得
\begin{align}
\langle\nu,q\rangle-\langle\nu^\pm,\lambda\rangle\langle\nu,\eta\rangle=\langle\nu^\pm,q-\lambda\eta\rangle.
\label{align:ch17-1765}
\end{align}

接着，对 Lax 对 $(\eta_{-k},q_{-k})$ 与 $(\eta_k,q_k)$ 应用\eqref{align:ch17-17310}，% 原书此处误排为 (16.3.10)
得到
\begin{align}
\frac{\langle\nu,q_k\rangle}{\langle\nu,\eta_k\rangle}-\frac{\langle\nu,q_{-k}\rangle}{\langle\nu,\eta_{-k}\rangle}=\frac{\langle\nu,\eta_{-k}q_k-\eta_kq_{-k}\rangle}{\langle\nu,\eta_{-k}\rangle\langle\nu,\eta_k\rangle}.
\label{align:ch17-1766}
\end{align}
由\eqref{align:ch17-1764}，当 $k\to\infty$ 时\eqref{align:ch17-1766}左端趋于 $\langle\nu^+,\lambda\rangle-\langle\nu^-,\lambda\rangle$。另一方面，右端趋于零，因为分子为 $O(k^{-1})$，而对大的 $k$，
\begin{align}
\langle\nu,\eta_{\pm k}\rangle\geqslant c\exp\bigl[\pm\tfrac{k}{2}(z^-+z^+)\bigr].
\label{align:ch17-1767}
\end{align}
因此
\begin{align}
\langle\nu^-,\lambda\rangle=\langle\nu^+,\lambda\rangle.
\label{align:ch17-1768}
\end{align}

将\eqref{align:ch17-1765}与\eqref{align:ch17-1768}结合，
\begin{align}
\langle\nu^-,q-\lambda\eta\rangle=\langle\nu^+,q-\lambda\eta\rangle.
\label{align:ch17-1769}
\end{align}

我们对 $(\eta,q)=(\eta_k,q_k)$ 应用\eqref{align:ch17-1769}。由\eqref{align:ch12-估计1}，对大的 $k$，
\begin{align}
\begin{cases}
\langle\nu^-,q_k-\lambda\eta_k\rangle\leqslant C\dfrac{1}{k}\exp(kz^-)\\
\langle\nu^+,q_k-\lambda\eta_k\rangle\geqslant c\dfrac{1}{k}\exp(kz^+),
\end{cases}
\label{align:ch17-17610}
\end{align}
这就对 $z^-<z^+$ 给出了预期的矛盾。类似地可证 $w^-=w^+$，于是 $\mathscr{R}$ 坍缩为单点。证明完成。
\end{proof}

使用上述定理为方程组\eqref{align:ch17-1731}构造解时的绊脚石在于：在撰写本书时，尚未确立任何现有格式所产生的近似解序列在 $L^\infty$ 中有界。因此，有界性不得不作为一个外加的（且令人厌烦的）假设。另一方面，一旦把有界性视为当然，验证\cref{theorem:ch17-两律真非线性补偿紧致}的另一要求\eqref{align:ch17-1761}并不困难。特别地，当序列 $U_k$ 由相对于抛物方程组\eqref{align:ch17-1733}的消失黏性方法生成时，条件\eqref{align:ch17-1761}直接由17.3节的讨论得出，因为两个守恒律的真非线性方程组总是配备一致凸熵。例如，如\cref{section:ch12-熵熵流对与速矢变换}所示，
在规范化条件\eqref{align:ch12-相互作用}下，
对充分大的 $k$，Lax 熵 $\eta_k$ 为凸。于是我们有

\begin{theorem}[两律方程组的消失黏性极限]\label{theorem:ch17-两律消失黏性}
对 $\mu>0$，设 $U_\mu$ 是两个守恒律的真非线性抛物方程组\eqref{align:ch17-1733}在上半平面上以\eqref{align:ch17-17311}为初值的解，其中当 $\mu\downarrow0$ 时 $U_{0\mu}\rightharpoonup U_0$ 在 $L^\infty(-\infty,\infty)$ 弱$^*$ 意义下成立。假设族 $\{U_\mu\}$ 含于 $L^\infty$ 的一个有界子集。则存在序列 $\{\mu_j\}$（$\mu_j\to0$，$j\to\infty$），使得 $\{U_{\mu_j}\}$ 在上半平面上几乎处处收敛于\eqref{align:ch17-1731}的一个解 $\bar U$，其初值为 $\bar U(x,0)=U_0(x)$（$-\infty<x<\infty$）。
\end{theorem}

对于由一类三点依赖域的单步差分格式生成的近似解序列，可得到完全类似的结果：
\begin{align}
U(x,t+\Delta t)-U(x,t)=\frac{\alpha}{2}G\bigl(U(x,t),U(x+\Delta x,t)\bigr)-\frac{\alpha}{2}G\bigl(U(x-\Delta x,t),U(x,t)\bigr),
\label{align:ch17-17611}
\end{align}
其中 $\alpha=\Delta t/\Delta x$ 是网格长度之比，而 $G$（可依赖于 $\alpha$）是满足相容性条件 $G(U,U)=F(U)$ 的函数。该类包括 Lax--Friedrichs 格式，其中
\begin{align}
G(V,W)=\frac{1}{2}[F(V)+F(W)]+\frac{1}{\alpha}(V-W),
\label{align:ch17-17612}
\end{align}
也包括 Godunov 格式，其中 $G(V,W)$ 表示以 $V$ 为左状态、$W$ 为右状态的\eqref{align:ch17-1731}Riemann 问题解的波后状态。在这些情形下，近似解在 $L^\infty$ 中的一致有界性条件同样须外加施加。

\end{document}
